% Generated draft: requires mathematical and source review.
\documentclass[11pt]{article}
\usepackage[margin=25mm]{geometry}
\usepackage{fontspec}
\setmainfont{Latin Modern Roman}
\usepackage{amsmath,amssymb,mathtools,unicode-math}
\setmathfont{Latin Modern Math}
\usepackage{enumitem,needspace,xurl,hyperref}
\setlength{\parindent}{0pt}
\setlength{\parskip}{5pt}
\setlength{\emergencystretch}{3em}
\newcommand{\N}{\mathbb{N}}
\newcommand{\Z}{\mathbb{Z}}
\newcommand{\Q}{\mathbb{Q}}
\newcommand{\R}{\mathbb{R}}
\newcommand{\C}{\mathbb{C}}
\newcommand{\F}{\mathbb{F}}
\newcommand{\A}{\mathbb{A}}
\newcommand{\PP}{\mathbb{P}}
\newcommand{\E}{\mathbb{E}}
\newcommand{\eps}{\varepsilon}
\newcommand{\dd}{\,\mathrm{d}}
\begin{document}
\section*{UnsolvedMath ID: 701146}
\textbf{Source problem code:} catalogue entry
\subsection{Definitions and mathematical statement}
A profinite group is a topological group isomorphic to an inverse limit of finite discrete groups. All homomorphisms below are required to be continuous. For profinite groups \(G\) and \(H\), their free profinite product \(G\amalg H\), denoted \(G*H\) in the source, is their coproduct in the category of profinite groups: it has canonical maps \(i_G:G\to G\amalg H\) and \(i_H:H\to G\amalg H\), and for every profinite group \(K\) and continuous homomorphisms \(f:G\to K\), \(g:H\to K\), there is a unique continuous homomorphism \(u:G\amalg H\to K\) satisfying \(u\circ i_G=f\) and \(u\circ i_H=g\).

For an integer \(n\geq 0\), let \(F_n\) be the abstract free group on \(n\) generators and let
\[
\widehat F_n=\varprojlim_{\substack{N\trianglelefteq F_n\\[F_n:N]<\infty}}F_n/N
\]
be the free profinite group of rank \(n\). In particular, \(\widehat F_0\) is the trivial group. All isomorphisms below are isomorphisms of topological groups.

The problem asks whether, for every pair of profinite groups \(G,H\),
\[
\bigl(\exists n\in\mathbb Z_{\geq0}\; G\amalg H\cong\widehat F_n\bigr)
\quad\Longrightarrow\quad
\bigl(\exists r,s\in\mathbb Z_{\geq0}\; G\cong\widehat F_r\ \text{and}\ H\cong\widehat F_s\bigr).
\]
This expresses the source question of whether both factors must be free profinite groups. The finite-rank formulation of the conclusion is equivalent here because the factors are continuous retracts of the product and hence are topologically finitely generated under the hypothesis. No nontriviality assumption on either factor is imposed.

The source adds, citing J. Neukirch, that the analogous assertion may fail when the free profinite product has infinite rank. This is retained as a supplied source qualification; the cited result and its precise scope have not been independently checked.
\subsection{Short English statement}
If the free profinite product of two profinite groups is a free profinite group of finite rank, must both factors be free profinite groups? The supplied record attributes the question to O. V. Mel'nikov and labels it open. It also cautions that the analogous assertion may fail in infinite rank.
\subsection{Sources}
\begin{itemize}
\item Local UnsolvedMath registry, supplied complete record with permanent ID \(701146\), exact source problem code \texttt{catalogue entry}, and title \emph{Kourovka Notebook Problem 15.58}. The supplied status is \texttt{open}; this label is not an independently verified literature assessment.
\item O. V. Mel'nikov, problem 15.58, \emph{The Kourovka Notebook (Unsolved Problems in Group Theory)}, 15th issue (2002). The supplied background also identifies the 21st issue edition (2026), \texttt{arXiv:1401.0300v48}. The collection metadata names E. I. Khukhro and V. D. Mazurov as editors.
\item J. Neukirch, \emph{Arch. Math.}, volume 22, no. 4 (1971), pp. 337--357. This reference is supplied for the qualification concerning infinite rank. No article title is provided, and the reference has not been independently consulted.
\end{itemize}
\end{document}
